\section{一审普通程序}
\begin{enumerate}
  \item 书记员确认双方当事人到场，审判长核对身份。
  \item 几个时间点。
        \begin{enumerate}
          \item 提交答辩状期间/答辩期。一（异）坨大便。答辩、提管辖权异议。
          \item 举证期限届满前。跟证据和证明有关。
          \item 法庭辩论终结前。请求、回避。诉讼请求要定稿、申请回避。
                \begin{enumerate}
                  \item 原告 \yjt 增加、变更请求。
                  \item 被告 \yjt 反诉。
                  \item 有独三 \yjt 有独三之诉。
                \end{enumerate}
          \item 判决宣告前。含“撤”字。
        \end{enumerate}
  \item 普通程序。体系最完整，内容最丰富的程序，是其他程序的范本模板。
  \item 起诉。
        \begin{enumerate}
          \item \bs{条件 = 立案条件}。考试绕着弯来……
                \begin{enumerate}
                  \item \zd{原告适格}。
                  \item \zd{有明确（到人，不到住所）的被告}。\al{被告要有诉讼权利能力 \dui。被告适格 \cuo}。
                  \item \zx{有具体的诉讼请求、事实、理由}。\al{原告应提供支持自己诉讼请求的证据 \cuo}。
                  \item \zd{属于法院主管 + 管辖}。
                \end{enumerate}
          \item 7 日内处理。只看条件，不看其它！！
                \begin{enumerate}
                  \item 符合条件 \yjt 登记立案 + 通知当事人。
                  \item 不符合 \yjt \yi \zx{裁定不予受理}。\er \zx{已受理 \yjt 裁定驳回起诉}。
                  \item 3 个会一起考的文书。裁定不予受理、裁定驳回起诉、判决驳回诉讼请求（败诉）。
                \end{enumerate}
          \item 起诉效果。仅一个时效中断。
          \item 起诉状。可口头、可书面。
                \begin{enumerate}
                  \item 内容就 4 样。当事人的有关情况、诉讼请求、事实和理由、证据相关。
                  \item 没有案由（案件进行分类）这个东西。
                \end{enumerate}
        \end{enumerate}
  \item 先行调解/人民调解看是否已立案。\yi 没立案 \yjt 确认调解协议效力。\er 立案后 \yjt 制作调解书。
  \item 受理。都可以自己推。看到驳回起诉 \yjt 考不予受理 \yjt 受理条件 \yjt 起诉条件 \yjt 当事人是否符合条件。
        \begin{enumerate}
          \item \yi 不属于管辖 \yjt 告知向有管辖权的法院起诉，仍坚持不予受理。\er 受理后发现不属于本院管辖 \yjt 移送。
          \item 存在有效的仲裁协议 \yjt 不予受理。受理后发现 \yjt 驳回起诉。
          \item 超过诉讼时效的处理。注意区分是否已被被判决驳回请求。
                \begin{enumerate}
                  \item 判决驳回请求，又起诉 \yjt 不予受理（重复起诉，一事不再理）。
                  \item 起诉 \yjt 受理（诉权未消灭）但驳回诉讼请求。\al{该反诉已过诉讼时效，应不予受理 \cuo}。
                \end{enumerate}
          \item 受理效果。不要和起诉效果混。
                \begin{enumerate}
                  \item 法院取得审判权。案件系属法院。有义务也有权利对案件作出裁判。
                  \item 取得排他管辖权。其他法院不得重复受理、当事人不得重复起诉。怪不得先到先得。
                  \item 当事人取得相对应的诉讼地位。据此行使诉讼权利，承担诉讼义务。
                \end{enumerate}
        \end{enumerate}
  \item \zx{一事不再理制度}。已经实体处理 \hz 正在实体审理重复起诉。害怕冲突判决。
        \begin{enumerate}
          \item 例外。撤诉不是啊！！没经过实体审理呢！！新的判断。重弱者。
                \begin{enumerate}
                  \item 未经实体处理。裁定不予受理/驳回起诉。撤诉/按撤诉处理（当事人同意）。
                  \item 3 费。抚养费、扶养费、赡养费。
                  \item 离婚没离成 \yjt 6 个月后、新情况、被告随时。
                  \item 维特收养关系的案件。新情况新理由。
                  \item 裁判生效后，发生新的事实，当事人再次提起诉讼的，法院应当受理。\al{没考虑后遗症、疫情}。
                \end{enumerate}
          \item \bs{最高法标准}。
                \begin{enumerate}
                  \item 后诉与前诉的\zd{当事人相同}。
                  \item 后诉与前诉的\zd{诉讼标的相同}。
                  \item 后诉与前诉的\zd{诉讼请求相同 \hz 后诉的诉讼请求在实质上否定前诉裁判结果}。
                \end{enumerate}
          \item 理论和实务很复杂的问题。民诉理论 3 个顶级难题。诉讼标的理论、既判力理论、诉讼系属理论。
        \end{enumerate}
  \item 一审必开庭，不分啥程序。开庭是为了查事实、查证据。
  \item 撤诉和缺席判决。都来 \yjt 对席判决。主动方不来 \yjt 视为撤诉。被动方不来 \yjt 缺席判决。
  \item \zd{诉讼障碍}。16 个情形，不要死记硬背。考试时候凭感觉选完全没问题。
        \begin{enumerate}
          \item 延期审理/延期开庭用决定决定决定。这次开不了，换个时间再开。
          \item 诉讼中止用裁定裁定裁定。要不要继续没谱。
          \item 诉讼终结。诉讼中死亡。诉讼中 = 受理后至判决生效。
                \begin{enumerate}
                  \item \zd{继受人放弃诉讼权利}、离婚诉讼中一方死亡。
                  \item 做题。离婚吗 \,? 终结 \,: [提到了别的人 \,? 中止（流程在谁谁管） \,: 终结]。
                  \item 离婚 vs 无效婚宣告制。
                        \begin{enumerate}
                          \item 无效事由要持续到起诉时，否则视为不存在。
                          \item 不允许调解、撤诉。
                          \item 申请时间。夫妻生前均可申请，夫妻一方 || 双方死亡 1 年内亦可。
                          \item 请求撤销婚姻 \bdy 请求解除婚姻关系。所以死了不是终结诉讼，而是继续审理。
                        \end{enumerate}
                \end{enumerate}
        \end{enumerate}
  \item 一审审期 6 + 6（院长批）。
  \item 判决。\zd{遗漏判决 \yjt 应当针对原告的诉讼请求作出判决}。
        \begin{enumerate}
          \item 一律公开宣判。告知上诉权。
          \item 离婚告知判决生效前不准另行结婚。
          \item 当庭宣判 10 日给，定期宣判立即给。当庭宣判 10 日给判决书，定期宣判当庭/立即给判决书。
          \item 生效判决实质错误，只能走再审，不能裁定补正。\al{笔误、计算、文字错误}。\al{报请上级批准后裁定补正 \cuo}。
          \item 裁判文书的公开。了解即可。只能查生效的不涉密的。
          \item 裁定、生效裁定、判决可以很多个，但生效判决就一个。
        \end{enumerate}
\end{enumerate}
